\documentclass[11pt,a4paper,openany]{ctexrep}
\usepackage[a4paper,margin=23mm,headheight=15pt]{geometry}
\usepackage{amsmath,amssymb,mathtools,bm}
\usepackage{siunitx}
\usepackage{booktabs,tabularx,array,longtable}
\usepackage{enumitem}
\usepackage{graphicx}
\usepackage{xcolor}
\usepackage{tikz}
\usetikzlibrary{arrows.meta,positioning,calc}
\usepackage[most]{tcolorbox}
\usepackage{hyperref}
\usepackage{fancyhdr}
\usepackage{titlesec}
\usepackage{microtype}
\usepackage{setspace}

\definecolor{KidBlue}{HTML}{22577A}
\definecolor{KidTeal}{HTML}{2A9D8F}
\definecolor{KidGold}{HTML}{E9C46A}
\definecolor{KidRed}{HTML}{C8553D}
\definecolor{KidGray}{HTML}{F3F5F7}
\definecolor{KidDark}{HTML}{263238}

\hypersetup{colorlinks=true,linkcolor=KidBlue,urlcolor=KidBlue,citecolor=KidTeal,pdftitle={KID 器件加工与实验入门},pdfauthor={KID Intro Guide project}}
\setstretch{1.15}
\setlength{\parindent}{2em}
\setlength{\parskip}{0.25em}
\setlist[itemize]{leftmargin=2em,itemsep=0.25em,topsep=0.4em}
\setlist[enumerate]{leftmargin=2em,itemsep=0.25em,topsep=0.4em}

\pagestyle{fancy}
\fancyhf{}
\fancyhead[L]{KID 器件加工与实验入门}
\fancyhead[R]{\leftmark}
\fancyfoot[C]{\thepage}
\renewcommand{\headrulewidth}{0.4pt}

\titleformat{\chapter}[display]{\bfseries\Huge\color{KidDark}}{\filleft\color{KidBlue}\fontsize{44}{44}\selectfont\thechapter}{1ex}{\titlerule\vspace{1ex}\filright}
\titleformat{\section}{\Large\bfseries\color{KidBlue}}{\thesection}{0.7em}{}
\titleformat{\subsection}{\large\bfseries\color{KidTeal}}{\thesubsection}{0.7em}{}

\newtcolorbox{chaptergoal}{colback=KidBlue!5,colframe=KidBlue,boxrule=0.8pt,arc=2mm,title=本章目标,fonttitle=\bfseries}
\newtcolorbox{processbox}{colback=KidTeal!6,colframe=KidTeal,boxrule=0.8pt,arc=2mm,title=工艺 → KID 因果链,fonttitle=\bfseries}
\newtcolorbox{measurebox}{colback=KidGold!16,colframe=KidGold!80!black,boxrule=0.8pt,arc=2mm,title=应该记录什么,fonttitle=\bfseries}
\newtcolorbox{warningbox}{colback=KidRed!5,colframe=KidRed,boxrule=0.8pt,arc=2mm,title=边界与安全,fonttitle=\bfseries}
\newtcolorbox{checkpointbox}{colback=KidGray,colframe=KidDark!55,boxrule=0.7pt,arc=2mm,title=理解检查,fonttitle=\bfseries}

\newcommand{\processchain}[1]{\begin{processbox}#1\end{processbox}}
\newcommand{\kidparam}[1]{\ensuremath{\mathrm{#1}}}

\begin{document}

\begin{titlepage}
\centering
\vspace*{2.0cm}
{\fontsize{30}{38}\selectfont\bfseries\color{KidDark} KID 器件加工与实验入门}\par
\vspace{0.5cm}
{\Large\color{KidBlue} 从版图到可测芯片}\par
\vspace{1.0cm}
{\large Volume 2 · v0.2 Development Draft · 2026-09-16}\par
\vfill
\begin{tcolorbox}[width=0.9\textwidth,colback=KidGray,colframe=KidBlue!55,arc=3mm]
\centering\large
本册不试图把读者训练成微纳工艺工程师。\\[0.4em]
它只做一件事：把每一道工艺步骤重新接回 KID 的可测参数，建立\\[0.4em]
\textbf{工艺 → 材料/几何 → 电磁参数 → 谐振器 → 低温测量}\\[0.4em]
这一条可诊断、可追溯、逐步可计算的工程链。
\end{tcolorbox}
\vfill
{\large 适合：第一次进入微纳平台前预习、加工记录、低温测试后回溯}\par
\end{titlepage}

\setcounter{tocdepth}{1}
\tableofcontents
\clearpage

\chapter{一块 KID 芯片是怎样“出生”的}

\begin{chaptergoal}
读完本章，你应该能把版图、薄膜、光刻、图形转移、切片封装、室温表征和低温测量串成一条完整流程；能解释为什么 GDS 不等于真实器件；并开始用 Batch ID、process traveler、witness sample 和 test coupon 为每一批器件保留可追溯证据。
\end{chaptergoal}

第一册从光子出发，最终走到 $S_{21}$。这一册把方向反过来：从一张 GDS 版图出发，追踪它怎样变成低温系统里的一条 resonance。

一个很容易产生的错觉是：\textbf{只要版图画对了，加工就是把图形“复印”到晶圆上。}真实情况不是这样。薄膜厚度会偏，材料参数会随位置变化；光刻会有 critical-dimension（CD）bias；刻蚀会改变边缘和侧壁；残胶、污染、切片和封装会继续改变器件。最终测到的 $f_0$、$Q_i$、$Q_c$ 和噪声，是整条制造链共同作用后的结果。

因此，加工学习的核心不是先背设备名和 recipe，而是建立一句始终不丢的主线：

\begin{processbox}
\centering
\textbf{工艺变量}
$\rightarrow$
\textbf{材料/真实几何}
$\rightarrow$
\textbf{电磁与超导参数}
$\rightarrow$
\textbf{谐振器可观测量}
$\rightarrow$
\textbf{阵列 yield 与科学可用性}
\end{processbox}

\section{先把“设计器件”和“真实器件”分开}

仿真中的器件通常由一组确定参数定义，例如：

\[
\{w,g,t,\epsilon_r,T_c,R_\square,\ldots\}_{\rm design}.
\]

加工之后，这些量变成具有偏差和分布的实际参数：

\[
\{w+\delta w,\;g+\delta g,\;t+\delta t,\;R_\square+\delta R_\square,\ldots\}_{\rm fabricated}.
\]

这里最重要的不是“误差总是存在”这句废话，而是进一步问三件事：

\begin{enumerate}
  \item 偏差是\textbf{系统性的}还是\textbf{随机的}？例如整片线宽都宽 $0.2\,\mu\mathrm{m}$，和少数位置随机缺口，处理方式完全不同。
  \item 偏差是\textbf{局部的}还是\textbf{wafer-scale 的}？膜厚从晶圆中心向边缘平滑变化，可能让整片 resonance map 发生有空间结构的扭曲。
  \item 偏差能不能\textbf{被测出来}？如果没有 CD、膜厚、方阻或见证片数据，低温以后只能猜。
\end{enumerate}

\section{一块单层 Al LEKID 的最小制造链}

作为第二册的主案例，我们先考虑结构相对清楚的单层 Al LEKID。不同平台的具体流程会不同，但逻辑上可以拆成下面几个阶段：

\begin{enumerate}
  \item \textbf{版图冻结：}确认 GDS/mask 版本、关键线宽/间距、测试结构和芯片编号。
  \item \textbf{基底与表面：}确认 substrate 批次、厚度、表面状态和预处理历史。
  \item \textbf{超导薄膜：}形成 Al 薄膜，并用见证片保留膜厚、$R_\square$、$T_c$ 等材料证据。
  \item \textbf{光刻：}把 mask/GDS 信息转移到 resist，产生真实的 CD 与边缘轮廓。
  \item \textbf{图形转移：}通过刻蚀或 lift-off 得到实际金属图形，并检查 residue、fence、过刻/欠刻等问题。
  \item \textbf{后处理与切片：}去胶、清洁、dicing，把 wafer 变成 die。
  \item \textbf{封装与 wire bond：}建立 feedline、ground 和 package 电磁边界。
  \item \textbf{室温/材料表征：}尽量在进入低温系统之前排除断线、短路、膜厚和明显 CD 问题。
  \item \textbf{第一轮 cooldown：}先做 resonance census 和基本 $f_0,Q_i,Q_c$ 统计，再进入更深的 optical/noise 测试。
\end{enumerate}

2014 年 McCarrick 等人的 150 GHz horn-coupled Al LEKID 原型就是一个很适合入门的参照：薄 Al 膜沉积在 Si wafer 上，再用标准光刻形成器件。它提醒我们，LEKID 可以拥有相对简洁的加工层数，但“层数少”并不等于“工艺变量不重要”。

\section{为什么每一步都必须留下证据}

建议从一开始就把信息分成三类。

\subsection{设计证据}

至少包括：

\begin{itemize}
  \item GDS / mask 版本和 commit/hash；
  \item resonator ID 与设计 $f_0$；
  \item meander linewidth/gap、IDC、coupler 等关键尺寸；
  \item test coupon / witness area 的位置；
  \item package 与 chip orientation 的设计版本。
\end{itemize}

\subsection{过程证据}

至少包括：

\begin{itemize}
  \item wafer/die/Batch ID；
  \item 每一步设备、SOP/recipe 标识和 run ID；
  \item 关键设定值与设备实际 readback；
  \item 时间、操作者、异常、返工；
  \item witness sample 与主器件的对应关系。
\end{itemize}

\subsection{测量证据}

至少逐步建立：

\begin{itemize}
  \item 膜厚与 wafer map；
  \item $R_\square$ 与可行时的 $T_c$；
  \item 关键 CD；
  \item 光学/SEM 等 inspection 图片；
  \item package/bond 图片；
  \item 低温 resonance map、$Q_i/Q_c$ 与后续 noise/optical 数据。
\end{itemize}

\begin{measurebox}
如果一个低温数据文件不能反查到具体 Batch ID、die ID、GDS 版本、薄膜 run 和封装版本，那么这批数据即使“看起来很好”，对下一版设计的工程价值也已经打了折扣。
\end{measurebox}

\section{Witness sample 和 test coupon 到底有什么用}

这两个词容易被混在一起。

\textbf{Witness sample（见证片）}通常跟随主样品经历某个材料或工艺步骤，用来做不方便直接在主芯片上完成的表征。例如同一次 Al deposition 中放一小块 witness sample，后续测膜厚、$R_\square$ 或 $T_c$。

\textbf{Test coupon / test structure（测试结构）}通常在版图中主动设计，用来测量某个加工或电学问题，例如 linewidth/gap monitor、连续性结构、不同耦合强度、不同 CD 的蛇形线等。

一个偏向“证明这一批材料是什么”，另一个偏向“证明这个工艺把几何做成了什么”。两者都不是浪费面积，而是把加工从黑箱变成可测系统。

\section{一个诊断例子：为什么谐振频率整体偏低？}

假设第一轮低温测试发现大部分 resonance 都存在，但 $f_0$ 普遍比设计值低。不要马上得出“电容画大了”这样的单因果结论。

因为
\[
f_0=\frac{1}{2\pi\sqrt{LC}},
\]
所以频率偏低至少意味着有效 $L$、有效 $C$ 或两者的组合比模型中更大。进一步可能来自：

\begin{itemize}
  \item 材料：$R_\square$、$T_c$ 或膜厚变化导致 $L_k$ 改变；
  \item 几何：meander / IDC 的 CD bias 改变 $L$ 或 $C$；
  \item 电磁环境：substrate / package 与仿真假设不同；
  \item 识别问题：测到的模式与设计 resonator 对应关系错误。
\end{itemize}

正确做法不是“选一个最像的原因”，而是优先调用已经保存的证据：膜厚和 $R_\square$ map、CD 抽样、resonance 的空间分布、package 版本。后面每一章都会把这棵诊断树继续拆细。

大阵列文献已经显示，这种思路不是洁癖：Shu 等人对 4-inch kilo-pixel LEKID 阵列的研究表明，实测 resonator 尺寸和 film thickness 可以解释相当一部分 resonance-frequency deviation；而 fabrication-induced frequency deviation 会进一步造成 frequency collision 并降低阵列 yield。

\section{本册的边界：学习“参数意义”，不是复制 recipe}

\begin{warningbox}
本册不会提供可以绕过设备培训的照抄 recipe。旋涂、曝光、显影、真空沉积、等离子体、湿法化学、高压和高温设备均应遵循所在平台 SOP 与设备管理员要求。这里要学的是：每个参数意味着什么、它改变什么、必须记录什么，以及异常如何传播到 KID 指标。
\end{warningbox}

\section{本章结束时，你应该能画出的图}

请自己画一条：

\[
\text{GDS}
\rightarrow
\text{film/process}
\rightarrow
\text{real geometry/material}
\rightarrow
\{L,C,R,Z_s\}
\rightarrow
\{f_0,Q_i,Q_c\}
\rightarrow
\text{yield/noise}.
\]

然后在每一根箭头旁边写出“我准备靠什么测量来验证它”。如果有一根箭头完全没有证据，这就是当前 fabrication plan 中最值得补的地方。

\begin{checkpointbox}
加工工程的目标不是让每个器件都与设计值一模一样，而是让偏差\textbf{可测、可解释、可预测、可反馈}。真正成熟的器件流程，最终会把“工艺误差”变成下一版仿真和版图的输入分布。
\end{checkpointbox}

\section*{进一步阅读}

\begin{itemize}
  \item H. McCarrick et al., ``Horn-Coupled, Commercially-Fabricated Aluminum Lumped-Element Kinetic Inductance Detectors for Millimeter Wavelengths,'' 2014, \href{https://arxiv.org/abs/1407.7749}{arXiv:1407.7749}.
  \item S. Shu et al., ``Understanding and minimizing resonance frequency deviations on a 4-inch kilo-pixel kinetic inductance detector array,'' 2021, \href{https://arxiv.org/abs/2105.14046}{arXiv:2105.14046}.
\end{itemize}

\chapter{洁净室、污染与工艺纪律}

\begin{chaptergoal}
理解洁净室不是“更干净的房间”，而是为了控制颗粒、有机物、金属交叉污染、湿度/静电和流程可重复性；建立 KID 加工所需的记录习惯。
\end{chaptergoal}

\section{先区分四类敌人}

对 KID 来说，“脏”并不是一个单一概念。至少应区分：颗粒导致断线或短路；有机残留改变界面与刻蚀；金属交叉污染改变材料性质；水汽和表面氧化改变界面状态。不同问题需要不同证据，不能把所有低 $Q_i$ 都笼统归因于“加工不干净”。

\processchain{污染/颗粒 $\rightarrow$ 薄膜或界面缺陷 $\rightarrow$ 局部电流拥挤/额外损耗/几何缺陷 $\rightarrow$ $Q_i$、yield、noise 异常}

\section{工艺纪律比单个 recipe 更重要}

第一次进入平台时，优先学会：样品如何编号、不同材料的工具是否分区、哪些设备允许哪些材料、样品如何暂存、何时允许用镊子/真空笔、怎样记录 recipe 版本、异常如何报备。重复性来自流程控制，而不是“记住上次旋钮调到多少”。

\begin{measurebox}
建议从第一批样品就建立 traveler：每一步都写下输入片号、输出片号、设备、recipe、开始/结束时间、异常、是否返工以及下一步允许条件。
\end{measurebox}

\chapter{基底与表面：器件真正开始的地方}

\begin{chaptergoal}
知道为什么 substrate 不只是机械支撑；理解介电常数、损耗、表面状态、平整度与清洗历史会进入谐振器和毫米波吸收问题。
\end{chaptergoal}

\section{选择基底时要问什么}

对微波谐振器，应关心介电常数、介质损耗、厚度及均匀性；对毫米/亚毫米波结构，还要考虑 substrate 中的传播、背短路或透镜几何。对加工，则要关心晶圆尺寸、平整度、表面粗糙度、是否有原生氧化层，以及后续沉积/刻蚀的兼容性。

\processchain{substrate / surface $\rightarrow$ 介电边界与界面状态 $\rightarrow$ 有效电容、场分布与损耗 $\rightarrow$ $f_0$、$Q_i$、optical absorption}

\section{“清洗”不是一个按钮}

清洗步骤应根据要去除的污染物和下一步工艺来定义。对于学习阶段，最重要的不是背化学配方，而是每次都明确：我要去掉什么？会不会留下新的表面状态？从清洗结束到沉积开始暴露了多久？这些问题应写进 process traveler。

\chapter{超导薄膜：从沉积参数到 $L_k$ 与 $Q_i$}

\begin{chaptergoal}
把“沉一层 Al”拆成可测的材料量：厚度 $t$、方阻 $R_\square$、$T_c$、均匀性、粗糙度、应力与界面；理解它们为什么会进入 KID 的谐振频率和损耗。
\end{chaptergoal}

\section{先学会区分设备名和物理量}

蒸发、溅射等是形成薄膜的手段；KID 真正“看到”的却是最后得到的材料状态。即使目标厚度相同，不同沉积历史也可能给出不同的 $R_\square$、$T_c$、晶粒/缺陷和界面状态。因此工艺记录不能只写“Al, 40 nm”，而应把 run 与后续表征绑定起来。

\section{第一条必须建立的量化链}

对足够薄、近似均匀的膜，可以先用
\[
R_\square = \frac{\rho}{t}
\]
建立最基本直觉：厚度和电阻率共同决定方阻。对超导薄膜，$R_\square$、$T_c$ 与材料模型又会进一步影响 sheet kinetic inductance。真正的定量关系要结合材料状态与适用近似，本册后续会把它与第一册的 $L_{k,\square}$ 接起来。

\processchain{$t,\rho,T_c,\text{disorder/interface}$ $\rightarrow$ $R_\square,L_{k,\square}$ 与损耗 $\rightarrow$ $L_k,Q_i$ $\rightarrow$ $f_0$, linewidth, responsivity}

\section{均匀性为什么对阵列格外重要}

单个谐振器的频率偏一点也许还能测；数百个频分复用谐振器如果出现相关的材料梯度，就可能让 frequency map 整体扭曲并增加 collision。阵列时代，wafer map 比“中心点膜厚是多少”更重要。

\begin{measurebox}
薄膜学习阶段至少建立：目标厚度与实际厚度、多个位置的均匀性、室温方阻、可行时的 $T_c$、沉积 base pressure / rate 等设备读回量、见证片编号，以及后续低温 resonance map。
\end{measurebox}

\chapter{光刻：GDS 并不等于最终几何}

\begin{chaptergoal}
理解 resist、曝光、显影和设备分辨率如何把版图变成带有 CD bias、圆角和轮廓的真实图形；学会从 KID 参数反推哪些几何最敏感。
\end{chaptergoal}

\section{最关键的概念：critical dimension}

对于 KID，meander linewidth、gap、IDC finger/gap、coupler 间隙往往比“大图形轮廓好不好看”更重要。设计阶段应先标出真正的 critical dimensions（CD），加工后再针对这些位置测量，而不是随机拍几张显微照片。

\processchain{曝光/显影偏差 $\rightarrow$ linewidth 与 gap 偏差 $\rightarrow$ $L_g,L_k,C,C_c$ 改变 $\rightarrow$ $f_0,Q_c$ 与吸收几何改变}

\section{为什么要测 bias}

如果版图上画 $2.0\,\mu\mathrm{m}$，实际稳定得到的是 $2.2\,\mu\mathrm{m}$，这不是一句“设备误差”就结束了。它应成为可量化的 process bias，并反馈到下一版 mask/GDS。DfM 的核心之一，就是把这类稳定偏差从“事故”变成“模型参数”。

\chapter{刻蚀与 Lift-off：把胶上的图形变成器件}

\begin{chaptergoal}
理解 subtractive etch 与 additive lift-off 的基本选择逻辑；知道侧壁、残胶、过刻、欠刻、redeposition 和金属毛刺为什么会成为电磁与可靠性问题。
\end{chaptergoal}

\section{两条常见路线}

一种路线是先沉整层薄膜，再用光刻胶定义保护区并刻蚀掉不需要的材料；另一种是先形成 resist opening，再沉积材料并 lift-off。二者都能得到图形，但对膜连续性、侧壁、最小间隙、材料兼容性和表面暴露历史的要求不同。

\processchain{图形转移偏差 $\rightarrow$ 实际边缘/侧壁/残留 $\rightarrow$ 电流路径与局部场改变 $\rightarrow$ $L,C,Q_i$, shorts/opens, yield}

\section{不要只看“有没有图形”}

低倍显微镜能确认大结构，但不能替代关键 CD、侧壁和 residue 的抽查。尤其是窄 gap、IDC、coupler 与 meander 转角，应提前定义接受标准和抽样位置。

\chapter{切片、封装与 Wire Bond：加工并没有在晶圆上结束}

\begin{chaptergoal}
理解 dicing、die handling、芯片固定、接地、wire bond、封装腔体、连接器和光学背腔都会改变器件的机械与电磁边界；建立 package revision 与 bond map 的记录习惯；学会区分“裸芯片本身的问题”和“装进盒子以后才出现的问题”。
\end{chaptergoal}

很多器件在 wafer 上是好的，装进 package 后却出现 baseline ripple、spurious resonance、$Q_i$ 下降或 feedline transmission 异常。原因并不神秘：\textbf{封装本身就是 microwave circuit 的一部分。}

\processchain{dicing / die attach / package / bond / connector $\rightarrow$ ground return、寄生 $L/C$、腔体与缝隙边界、光学位置 $\rightarrow$ feedline / resonator 环境改变 $\rightarrow$ $S_{21}$ baseline、spurious modes、$Q_i/Q_c$、串扰与 optical coupling}

\section{Dicing：从 wafer 变成 die 的第一道风险}

切片不是“把晶圆切小”这么简单。对 KID 来说，需要关心至少四类后果：

\begin{itemize}
  \item edge chipping / crack 是否侵入有效器件区；
  \item 颗粒、切削残留或保护材料是否污染金属和 gap；
  \item die 尺寸公差是否影响 package fit、背短路间距或 optical alignment；
  \item die handling 是否增加划伤、静电或局部机械应力风险。
\end{itemize}

因此 GDS 阶段就应定义 dicing street / keep-out，而不是 wafer 做完后再临时寻找下刀位置。

\begin{warningbox}
实际 dicing 参数、保护层、清洗和机械操作必须由目标平台 SOP 与设备管理员决定。本章不提供可直接执行的 blade / speed / coolant 等 recipe，只讨论 KID 需要验收哪些结果。
\end{warningbox}

\section{Die attach：固定方式也会进入电磁边界}

芯片需要被固定、热化并保持目标位置。任何 adhesive、clamp、metal contact 或 backside support 都可能同时影响：

\begin{itemize}
  \item thermal path：芯片是否可靠热化到冷台；
  \item mechanical stress：降温收缩后是否产生应力；
  \item microwave boundary：substrate backside 周围是否出现额外电场/电流路径；
  \item optical geometry：背短路、horn、lens 或 absorber-to-backshort spacing 是否保持设计值。
\end{itemize}

所以“同一片芯片换一个盒子”并不是严格意义上的同一实验。如果 package geometry 或 die attach 改了，应当作为新的 package revision 记录。

\section{Wire bond 不只是直流连接}

一根 bond wire 在 GHz 下具有明显电感，它不是理想零阻抗短路。大量 ground bonds 的作用也不只是“接地”，而是帮助芯片 ground 与 package ground 在微波频段保持接近等势，并抑制某些 slotline / common-mode 激励。

如果 ground plane 两侧被 feedline gap 分开，却没有足够的跨接或 package return path，就可能出现设计中没有的奇模或寄生传播路径。

已有 superconducting microwave chip-mount 研究直接展示了 wirebond arrangement、chip-to-box capacitance 与 cavity / slotline / crosstalk mode 之间的关系。这说明在 KID module 中，bond map 应该像 GDS 一样有版本，而不是“现场看着多打几根”。

\subsection{Signal bond 与 ground bond 要分开思考}

signal bond 需要保持 feedline transition 的阻抗连续性；ground bond 更关心回流路径和地平面等势。两类 bond 的失败模式不同：

\begin{itemize}
  \item signal bond 异常更容易直接表现为 transmission 很差或断路；
  \item ground return 不好可能仍然“有信号”，但 baseline、spurious modes、crosstalk 或 resonance line shape 会变差。
\end{itemize}

\section{Package cavity 与 seam：盒子也会谐振}

任何有限尺寸的导电腔体都存在 electromagnetic modes。package、lid、chip、PCB、connector 和空隙共同决定这些 mode 的频率与耦合强弱。

如果 package mode 靠近 KID readout band，它可能表现为：

\begin{itemize}
  \item 很宽或很深、与单个 resonator 不同的 $S_{21}$ 结构；
  \item 大范围 baseline ripple；
  \item 多个 resonator 同时受到影响；
  \item 换 package/lid/bond 后频率或形状明显变化。
\end{itemize}

更微妙的是 package conductor / seam loss：即使没有明显 spurious resonance，电磁场或感应电流如果参与了有损 package 区域，也可能限制可达到的 $Q_i$。近年来的 superconducting resonator experiments 已经能通过专门几何把 surface、bulk 和 package loss 分离出来。

因此，看到 $Q_i$ 不高时，不应自动把责任全部推给超导膜。

\section{150 GHz LEKID 还多了一层：光学封装}

对于 horn-coupled 或 backshort-coupled LEKID，module 同时有两套电磁问题：

\begin{enumerate}
  \item GHz readout：feedline、bond、ground、connector、package modes；
  \item 150 GHz optical path：waveguide/horn、substrate、absorber positioning、backshort 与缝隙。
\end{enumerate}

这两套系统可能共用同一块金属盒，但容差关注点完全不同。某个 mechanical spacer 在 GHz 看起来只是结构件，在 150 GHz 可能已经是重要的 electrical length。

因此 module drawing 中应明确区分：

\begin{itemize}
  \item microwave critical dimensions；
  \item optical critical dimensions；
  \item purely mechanical clearances。
\end{itemize}

\section{Package revision 必须和 chip revision 解耦}

建议从第一批器件就使用两套版本号：

\begin{center}
\texttt{Chip GDS Rev C03} \qquad + \qquad \texttt{Package Rev P02}
\end{center}

一次 cooldown 的 sample identity 至少应能追溯到：

\[
\text{wafer/die ID}
+\text{GDS rev}
+\text{film batch}
+\text{package rev}
+\text{bond map rev}.
\]

否则半年后看到“C03 的 $Q_i$ 很差”，你可能根本不知道那次测试用的是哪个 package。

\section{建立 Bond Map，而不是只拍一张照片}

bond map 可以非常简单：在 package drawing 或高分辨率装配照片上标出每根或每组 bond 的功能和位置，例如 signal-in、signal-out、ground stitching、跨 feedline ground bridge 等。

返工时也要更新。因为“断了一根又补了两根”对 GHz 电路并不是完全无关的小修。

\begin{measurebox}
\textbf{Package / Bond Card 最小字段：}
package revision；die ID 与 orientation；die attach 方法/材料 ID；关键 spacer/backshort dimension；connector / PCB / transition revision；signal bond 数量与位置；ground bond map；装配照片；返工记录；continuity test；对应 cooldown ID。
\end{measurebox}

\section{室温先做哪些 package sanity check}

在送进低温系统前，可以先做不昂贵的检查来避免浪费一次 cooldown：

\begin{itemize}
  \item die 是否坐平、方向是否正确；
  \item 芯片边缘是否有明显崩裂或压伤；
  \item wire bond 是否有明显脱落、互碰、跨错 pad；
  \item signal path continuity 是否合理；
  \item connector / transition 是否固定；
  \item optical stack 的关键 spacer / lid / backshort 是否装对 revision；
  \item 拍摄可追溯的最终装配照片。
\end{itemize}

对可在室温进行的 VNA / continuity 预检，应把它当作“找灾难性错误”的 gate，而不是期待看到低温 superconducting resonance。

\section{如何区分 chip problem 和 package problem}

一个很有用的思路是做 package swap 或 reference structure。

如果同一 die 在不同 package 中表现显著不同，或者多个不同 die 在同一个 package 中都出现相似 baseline/spurious features，那么 package hypothesis 的优先级会上升。

反过来，如果异常随着 die / wafer 走，而 package 保持不变，则更应回查 film / lithography / etch。

这仍然不是数学证明，但它是一种非常便宜的 controlled experiment。

\begin{checkpointbox}
当你看到一条难看的 $S_{21}$ baseline 时，现在应该先问：这是单个 resonator 的问题，还是整个 feedline/package 的问题？换 lid、bond、connector 或 package 后是否跟着变？如果没有 package revision 和 bond map，这个问题往往根本无法回答。
\end{checkpointbox}

\section*{进一步阅读}

\begin{itemize}
  \item J. Wenner et al., wirebond crosstalk and cavity/slotline modes in superconducting chip mounts, \href{https://arxiv.org/abs/1011.4982}{arXiv:1011.4982}.
  \item Y. Huang et al., conducting package losses limiting superconducting resonator $Q_i$, \href{https://arxiv.org/abs/2304.08629}{arXiv:2304.08629}.
  \item \textit{Surpassing millisecond coherence in on-chip superconducting quantum memories...}, multimode separation of surface, bulk and package losses, \href{https://www.nature.com/articles/s41467-024-47857-6}{Nature Communications 15, 2024}.
\end{itemize}

\chapter{工艺表征：在进冰箱之前尽量多知道一点}

\begin{chaptergoal}
建立“测什么—为什么测—它能排除什么”的表征思路，而不是把 profilometer、SEM、四探针等设备当成独立知识点。
\end{chaptergoal}

\section{常见表征与问题对应}

\begin{center}
\begin{tabularx}{\textwidth}{p{3.2cm} X X}
\toprule
手段 & 主要得到什么 & 对 KID 的用途 \\
\midrule
光学显微镜 & 大尺度图形、断线/短路、颗粒 & 快速 yield 筛查 \\
轮廓/膜厚测量 & 膜厚、台阶高度 & $t$、etch depth、run 均匀性 \\
四探针 & 方阻 $R_\square$ & 材料一致性、$L_k$ 先验 \\
SEM & CD、边缘、侧壁细节 & linewidth/gap bias、缺陷定位 \\
AFM（如可用） & 粗糙度、局部形貌 & 表面/界面质量证据 \\
$T_c$ 测量 & 超导转变 & 材料批次与模型输入 \\
\bottomrule
\end{tabularx}
\end{center}

\section{表征不是越多越好}

真正高价值的是能关闭假设空间的测量。例如阵列 $f_0$ 系统性偏移时，wafer-map 方阻/膜厚和关键 CD 往往比再拍一张漂亮的器件全景图更有诊断价值。

\chapter{低温测量如何反推工艺}

\begin{chaptergoal}
把低温 $S_{21}$ 从“最终成绩”变成 fabrication diagnosis：第一轮 cooldown 先建立 resonance inventory 和 wafer map，再看 $f_0,Q_r,Q_i,Q_c$、功率与温度依赖，最后才进入噪声与 NEP；学会用多条证据缩小 design / material / fabrication / package / readout 的责任范围，而不是把一个异常 notch 直接归因到某道工艺。
\end{chaptergoal}

第一次把一片新 KID wafer 放进低温系统时，很容易兴奋地直接去找最好的 $Q_i$ 或最低的 NEP。但从 fabrication learning 的角度，更重要的问题其实更朴素：\textbf{我们做出来的东西，整体上和设计有多像？偏差有没有结构？}

\processchain{低温 observables $\rightarrow$ resonance inventory / map $\rightarrow$ 与 film、CD、defect、package 数据对照 $\rightarrow$ 缩小 design / material / fabrication / package / readout 的 hypothesis space}

\section{第一轮 cooldown 的五个优先级}

建议按下面顺序，而不是一上来就做极限 sensitivity：

\begin{enumerate}
  \item 找到多少个 resonance？
  \item $f_0$ 相对设计整体偏高还是偏低？spread 多大？
  \item $Q_c$ 是否大致符合 coupling design？
  \item $Q_i$ 的分布怎样？是否有空间相关性或异常簇？
  \item baseline / spurious modes / package features 是否干扰正常拟合？
\end{enumerate}

只有这五项基本稳定之后，power sweep、temperature sweep、noise PSD、optical response 才更容易解释。

\section{先复习三个品质因数}

对一个常见 hanger resonator，可以写

\[
\frac{1}{Q_r}=\frac{1}{Q_i}+\frac{1}{Q_c},
\]

其中 $Q_r$（也常写 loaded $Q$）决定总线宽，$Q_i$ 表示内部损耗，$Q_c$ 表示与 feedline 的外部耦合。

这三个量在 fabrication diagnosis 中承担不同角色：

\begin{itemize}
  \item $f_0$：对 $L,C$ 以及材料/几何变化敏感；
  \item $Q_c$：优先反映 coupling geometry 与局部 EM environment；
  \item $Q_i$：综合承载 film、interface、TLS、quasiparticle、vortex、radiation、package 等内部/寄生损耗。
\end{itemize}

因此 $Q_i$ 最有价值、也最不能“看到一个数就下结论”。

\section{Step 1：建立 resonance inventory}

第一份低温输出应该不是漂亮的单个 notch，而是一张全 readout band 的 inventory：

\begin{itemize}
  \item design resonance count；
  \item detected resonance count；
  \item unmatched / ambiguous resonance count；
  \item frequency collision / merged-notch candidate；
  \item obvious spurious / package-mode candidate；
  \item unusable frequency regions。
\end{itemize}

定义一个最朴素的 detection yield：

\[
Y_{\mathrm{det}}=\frac{N_{\mathrm{matched}}}{N_{\mathrm{design}}}.
\]

它不等于 science yield，但能让不同 fabrication batch 有一个最基本的比较起点。

\section{Step 2：把 frequency 变成 map，而不是一列数字}

如果能把 measured resonance 对应回 physical resonator，就可以定义

\[
\Delta f_i=f_{0,i}^{\mathrm{meas}}-f_{0,i}^{\mathrm{design}},
\qquad
\epsilon_i=\frac{\Delta f_i}{f_{0,i}^{\mathrm{design}}}.
\]

然后至少看三种图：

\begin{enumerate}
  \item measured vs designed frequency；
  \item fractional error histogram；
  \item error vs physical position / wafer coordinate。
\end{enumerate}

第三张图特别重要，因为 fabrication variation 往往具有空间结构。已有大规模 MKID 工作正是通过 physical-pixel-to-frequency mapping 和 wafer-scale statistics 来理解 frequency scatter，并进一步做 trimming 或优化 resonator ordering。

\section{Frequency map 怎样反推 fabrication}

不同图形特征对应不同的优先假设：

\begin{itemize}
  \item \textbf{全片近似同方向整体偏移：} 先检查材料模型、平均 $R_\square/t/T_c$、systematic CD bias、substrate dielectric / EM model；
  \item \textbf{中心到边缘平滑梯度：} film thickness / $R_\square$ map、lithography CD map、etch nonuniformity；
  \item \textbf{局部区域异常：} defect map、颗粒、局部 lithography / etch failure；
  \item \textbf{相邻设计频率大量碰撞：} scatter 是否已经接近或超过设计 spacing；
  \item \textbf{随机但很大的 scatter：} 先确认 resonator mapping 和 fitting 是否可靠，再讨论 fabrication randomness。
\end{itemize}

不要一看到 frequency 偏差就自动修改 IDC。那可能只是在用设计补偿未知 fabrication bias。

\section{Step 3：看外部耦合品质因数，验证 coupling geometry}

如果 $Q_i$ 足够高，$Q_c$ 与设计偏差可以成为很有价值的 coupling-region check。

当一批器件的 $Q_c$ 系统性偏大或偏小时，应优先对照：

\begin{itemize}
  \item coupler gap / length final CD；
  \item feedline geometry；
  \item substrate / package environment 是否与 simulation 一致；
  \item resonance fit 是否正确处理 cable delay / asymmetry / background。
\end{itemize}

如果只有个别器件异常，则再查局部 lithography / bridging / defect。

\section{Step 4：看内部品质因数，但先把它当“症状”}

$Q_i$ 低可能来自很多不同来源：

\begin{itemize}
  \item film intrinsic loss / quasiparticles；
  \item substrate / interface TLS；
  \item etch damage / residue；
  \item trapped magnetic vortex；
  \item radiation / leakage；
  \item package conductor / seam；
  \item local defect；
  \item fitting / calibration error。
\end{itemize}

所以一个低 $Q_i$ 数值不构成 root cause。真正有用的是它随功率、温度、位置、geometry 或 package revision 怎样变化。

\section{Power sweep：区分“会饱和的损耗”和其他损耗}

很多 amorphous-interface TLS loss 会表现出明显的 microwave-power dependence：低激励时损耗较强，随着 TLS 被逐渐饱和，$Q_i$ 可能上升。

因此 power sweep 是检查 TLS-like behavior 的重要工具之一。但仍要注意：readout power 过高时，KID 自身会出现非线性、加热或 bifurcation 等效应。

所以 power sweep 的目标不是“功率越大 $Q$ 越好”，而是找到线性/低功率区、识别趋势，并与 geometry / surface-process split 比较。

\section{Temperature sweep：把材料模型重新接回来}

随着温度接近 $T_c$ 的一定比例，thermal quasiparticle population 增加，complex conductivity、kinetic inductance 和 dissipation 都会变化，因此 $f_0(T)$ 和 $Q_i(T)$ 可以与 Mattis--Bardeen / quasiparticle model 比较。

如果 measured $T_c$、$R_\square$、film thickness 与低温 $f_0(T)$ / $Q_i(T)$ 能形成一致解释，那么第 3 章的 film batch card 才真正闭环。

反过来，如果 room-temperature film proxy 很正常但低温行为异常，就要扩大到 interface、magnetic environment、package 等其他 hypothesis。

\section{Package / spurious mode 怎样从传输曲线看出来}

单个 KID resonance 通常是相对局域、可用 resonator model 拟合的 notch；package / cavity / feedline features 往往更宽、影响更大频段，或者同时扭曲多个 resonator。

可以通过几种便宜的比较提高判断力：

\begin{itemize}
  \item empty / reference package sweep；
  \item 同一 die 换 package / lid / bond revision；
  \item 多个不同 die 是否都在相似频率出现同一 feature；
  \item EM simulation 中的 package mode 是否与实测 feature 接近。
\end{itemize}

如果 package feature 没有先处理，强行拟合附近 KID 的 $Q_i/Q_c$ 往往没有太大意义。

\section{建立一个 fabrication diagnosis matrix}

\begin{center}
\begin{tabularx}{0.98\textwidth}{>{\bfseries}p{3.3cm} X X}
\toprule
低温症状 & 优先对照的数据 & 还不能排除 \\
\midrule
$f_0$ 全片整体偏移 & $R_\square,t,T_c$、systematic CD、EM model & substrate / calibration / model error \\
$f_0$ 有空间梯度 & film map、CD map、etch map & thermal / mapping error \\
大量 collision & frequency scatter statistics、spacing design & fitting/mapping ambiguity \\
$Q_c$ 系统偏差 & coupler CD、feedline、EM model & background-fit / package \\
$Q_i$ 低且 power-dependent & surface/interface split、geometry participation & heating / readout nonlinearity \\
$Q_i$ 低且随 package 变化 & bond/package/seam/ground revision & chip 与 package 共同损耗 \\
局部 resonator 缺失 & defect map、open/short、local CD & collision / 极端 frequency shift \\
宽带 baseline / spurious & package / connector / reference sweep & readout chain calibration \\
\bottomrule
\end{tabularx}
\end{center}

这张表的意义不是自动诊断，而是防止“一个症状对应一个原因”的思维捷径。

\section{从 resonance index 到 physical pixel：阵列必须做 mapping}

当 resonator 数量增加，仅知道“第 37 个 notch”没有意义。必须尽可能建立

\[
\text{frequency ID} \leftrightarrow \text{physical pixel ID}.
\]

大规模 MKID 阵列已经发展出 cryogenic LED mapper 等方法，把 illumination response 与 resonance frequency 对应起来；也有工作在完成 mapping 后再做 lithographic frequency trimming，提高 spacing uniformity 和 collision yield。

对我们的第一批小阵列，不一定需要复杂 LED mapper，但从数据结构一开始就应该预留 physical ID、GDS coordinate、design frequency 与 measured frequency 的映射表。

\section{什么时候才进入 noise / NEP}

如果下面这些问题还没解决：

\begin{itemize}
  \item resonance mapping 不清楚；
  \item $Q_c$ 与设计严重不符；
  \item package baseline 不稳定；
  \item readout power 已进入明显非线性；
  \item thermalization / temperature 记录不可靠；
\end{itemize}

那么直接比较 NEP 往往会把 fabrication、readout 和 calibration 混在一起。

因此第二册的顺序有意把 NEP 放晚：\textbf{先证明器件是谁、在哪、参数怎样，再讨论它有多灵敏。}

\begin{measurebox}
\textbf{First Cooldown 最小交付物：}
完整 raw $S_{21}$ sweep；resonance inventory；design-to-measured mapping；$f_0,Q_r,Q_i,Q_c$ 表；frequency-error histogram + spatial map；selected power sweep；selected temperature sweep；package/spurious-mode 标记；与 film/CD/defect/package metadata 的统一 sample ID。
\end{measurebox}

\section{最重要的输出不是“最高内部品质因数”，而是下一批该改什么}

一次 fabrication-learning cooldown 的最终报告最好只形成少量明确结论：

\begin{enumerate}
  \item 哪些 design assumptions 被验证？
  \item 哪些 fabrication variables 与异常有相关证据？
  \item 哪些 hypothesis 还没有区分？
  \item 下一批只改哪一两个变量最有信息量？
\end{enumerate}

如果一次 cooldown 最后变成“这片不太好，下一版都改改”，那么最珍贵的学习机会基本被浪费了。

\begin{checkpointbox}
看到一个低 $Q_i$ resonator 时，不要问“是哪道工艺做坏了？”先问：它在 wafer 上哪里？附近器件也低吗？power dependence 怎样？temperature dependence 怎样？package revision 是什么？film/CD/defect map 是否同位置异常？当这些信息齐了，root-cause analysis 才真正开始。
\end{checkpointbox}

\section*{进一步阅读}

\begin{itemize}
  \item C. R. H. McRae et al., resonator loss mechanisms and measurement methodology, \href{https://www.nist.gov/publications/materials-loss-measurements-using-superconducting-microwave-resonators}{NIST / Rev. Sci. Instrum. 91, 091101}.
  \item J. Gao et al., cryogenic LED pixel-to-frequency mapper for MKID arrays, \href{https://www.nist.gov/publications/cryogenic-led-pixel-frequency-mapper-kinetic-inductance-detector-arrays}{NIST / Appl. Phys. Lett. 111, 2017}.
  \item X. Liu et al., measured-frequency mapping followed by lithographic trimming, \href{https://www.nist.gov/publications/superconducting-micro-resonator-arrays-ideal-frequency-spacing}{NIST / Opt. Express 25, 2017}.
  \item C. M. McKenney et al., wafer-scale frequency-spacing statistics and collision-aware array layout, \href{https://www.nist.gov/publications/tile-and-trim-micro-resonator-array-fabrication-optimized-high-multiplexing-factors}{NIST / Rev. Sci. Instrum. 90, 023908}.
\end{itemize}

\chapter{Design for Fabrication：不要只设计一个“仿真里能工作的 KID”}

\begin{chaptergoal}
把 fabrication tolerance、process bias、test structures 和 array yield 纳入设计；让下一版 GDS 能主动吸收加工平台的真实能力。
\end{chaptergoal}

\section{从 nominal design 走向 process window}

工程设计不只问 nominal geometry 下性能是否最好，还要问在 $\pm\Delta w$、$\pm\Delta g$、$\pm\Delta t$、材料批次变化和封装偏差下，性能是否仍然可接受。对大阵列，frequency collision 往往让这种 tolerance 比单像素峰值性能更重要。

\processchain{工艺能力分布 $\rightarrow$ geometry/material distribution $\rightarrow$ $f_0,Q_i,Q_c$ distribution $\rightarrow$ collision / yield $\rightarrow$ mask 与版图修正}

\section{建议从第一版就放入的 test structures}

可以考虑独立的 linewidth/gap monitor、方阻/连续性结构、不同 coupler 强度、不同 meander CD、不同 IDC bias，以及用于测量 etch/膜厚的见证区域。具体结构应服务于“下一次设计要修正什么”，而不是为了填满版图空白。

\section{第二册最终要形成的闭环}

\begin{center}
\textbf{设计假设 $\rightarrow$ 工艺记录 $\rightarrow$ 室温表征 $\rightarrow$ 低温结果 $\rightarrow$ 参数反演 $\rightarrow$ 下一版设计}
\end{center}

当这个闭环可以重复运行时，加工就不再是外包给洁净室的黑箱，而成为 KID 设计模型的一部分。

\chapter{把工艺数据变成下一版设计：一个可运行的定量闭环}

\begin{chaptergoal}
把前面章节里的 film map、CD map、resonance map 和 DfM 连接成一个可以真正执行的数据流程：先定义可追溯的数据表，再建立一个透明的小信号 surrogate，检查它能解释多少 frequency scatter，最后只对可重复的 systematic component 做下一版设计补偿；同时理解为什么“相关”不等于“工艺原因已经证明”。
\end{chaptergoal}

到 v0.1 为止，我们已经知道应该测什么、记录什么、怎样避免把一个异常 notch 直接归因到某一道工艺。v0.2 往前再走一步：\textbf{把这些记录变成能跑的模型。}

这里故意先使用一个很小的 synthetic example，而不是一上来做复杂机器学习。原因很简单：如果一个 64-resonator toy model 都不能把 ID、单位、sensitivity、collision rule 和 correction 写清楚，那么换成几千个像素只会把混乱放大。

\processchain{wafer metrology map $\rightarrow$ per-resonator predictor $\rightarrow$ measured frequency residual $\rightarrow$ collision audit $\rightarrow$ repeatability check $\rightarrow$ next-GDS correction}

\section{第一步：所有数据必须先通过 resonator ID 对齐}

对第 $i$ 个 resonator，至少需要把下面几类量放到同一个可连接的数据系统中：

\begin{itemize}
  \item 设计侧：$f_{i}^{\mathrm{target}}$、GDS revision、物理坐标；
  \item 薄膜侧：$R_{\square,i}$、厚度或其他可空间插值的 material proxy；
  \item 几何侧：linewidth、gap、coupler 等 measured CD；
  \item 低温侧：$f_{0,i}^{\mathrm{meas}}$、$Q_{r,i}$、$Q_{i,i}$、$Q_{c,i}$；
  \item 追溯侧：batch / wafer / die / package / cooldown ID。
\end{itemize}

因此 v0.2 的第一件事不是“拟合一个模型”，而是先让
\[
\text{wafer coordinate}
\leftrightarrow
\text{physical pixel}
\leftrightarrow
\text{resonator ID}
\leftrightarrow
\text{measured notch}
\]
能够稳定连接。

本仓库从 v0.2 开始用 \texttt{notes/DATA\_CONTRACT.md} 固定这一层。对应的两个核心表是：

\begin{itemize}
  \item \texttt{wafer\_metrology\_map.csv}：保存物理位置对应的 film / CD / defect 数据；
  \item \texttt{resonator\_map.csv}：连接 design frequency、physical pixel 与低温 fit 结果。
\end{itemize}

\section{第二步：先定义我们到底在拟合什么}

最直接的目标量是 fractional frequency error：
\[
\epsilon_i
=
\frac{
f_{0,i}^{\mathrm{meas}}-f_{i}^{\mathrm{target}}
}{
f_{i}^{\mathrm{target}}
}.
\]

如果只看 $\epsilon_i$ 的 histogram，我们只能知道 spread 多大；如果把它与同位置的 film / geometry 数据对齐，就可以尝试建立 predictor。

对薄膜部分，第 3 章已经给出近似关系
\[
\frac{\delta f_0}{f_0}
\approx
-\frac{\alpha}{2}
\frac{\delta L_k}{L_k}.
\]

在一个只用于建立工程直觉的近似里，如果 $T_c$ 变化暂时忽略，并把 sheet kinetic inductance 的变化主要接到 $R_\square$，可以写成
\[
\epsilon_{R_\square,i}
\approx
-\frac{\alpha}{2}
\frac{R_{\square,i}-R_{\square,0}}{R_{\square,0}}.
\]

这部分至少有明确的物理方向。

几何部分则更适合写成 sensitivity：
\[
s_w
=
\frac{\partial \ln f_0}{\partial \ln w},
\qquad
\epsilon_{w,i}
\approx
s_w
\frac{w_i-w_0}{w_0}.
\]

这里的 $s_w$ \textbf{不是材料常数，也不是本讲义可以替实验室替你规定的数字}。它应该来自 Sonnet/CST/解析 LC sensitivity，或者来自已经验证过的 measured-data regression。

于是一个最小 predictor 可以写成
\[
\hat{\epsilon}_i
=
-\frac{\alpha}{2}
\frac{R_{\square,i}-R_{\square,0}}{R_{\square,0}}
+
s_w
\frac{w_i-w_0}{w_0}.
\]

真正值得看的不是 $\hat{\epsilon}_i$ 本身，而是 residual：
\[
r_i
=
\epsilon_i-\hat{\epsilon}_i.
\]

如果 residual 仍然和 wafer position、gap、package revision 或某个 process split 强相关，就说明模型还漏掉了结构。

\section{第三步：用一个 synthetic 64-resonator 数据集先把逻辑跑通}

仓库中的
\[
\texttt{examples/wafer\_to\_resonance\_demo.py}
\]
使用一个固定随机种子的 $8\times8$ synthetic array。它故意包含：

\begin{itemize}
  \item 一个带有空间梯度的 $R_\square$ map；
  \item 一个带有空间梯度的 linewidth map；
  \item 一个额外的小随机 residual；
  \item 64 个 nominally 等间隔的 design frequencies；
  \item 一个显式的 loaded-$Q$ 和 collision rule。
\end{itemize}

这个例子中的数值全部是教学数据，不代表任何真实 Al LEKID 工艺能力，也不应该被复制成 fabrication specification。

运行：

\begin{verbatim}
cd volume2-fabrication
python3 examples/wafer_to_resonance_demo.py
\end{verbatim}

脚本会生成：

\begin{itemize}
  \item \texttt{wafer\_metrology\_map.csv}；
  \item \texttt{resonator\_map.csv}；
  \item \texttt{summary.txt}。
\end{itemize}

固定 seed 下，这个 toy model 在补偿前得到大约
\[
\sigma_{\epsilon}\approx 860~\mathrm{ppm},
\]
模型补偿后剩余大约
\[
\sigma_{r}\approx 111~\mathrm{ppm}.
\]

同一个例子中，在“相邻实测 resonance 间隔小于 5 个 loaded linewidth”这个\textbf{仅用于演示}的判据下，collision pair 从 2 组降到 0 组。

\begin{warningbox}
这里最重要的不是 ``860 ppm''、``111 ppm'' 或 ``5 linewidths'' 这些数字。它们只是 deterministic regression target，方便 CI 检查示例没有被改坏。真正实验中，sensitivity、process distribution 和 collision margin 都必须重新定义。
\end{warningbox}

\section{第四步：为什么可以做 pre-compensation}

如果上一批已经得到一个可重复的 predictor $\hat{\epsilon}_i$，而下一批仍使用足够相似的 process，可以把目标频率 $f_i^{\mathrm{target}}$ 反推成新的 drawn/design frequency：

\[
f_i^{\mathrm{draw,new}}
=
\frac{
f_i^{\mathrm{target}}
}{
1+\hat{\epsilon}_i
}.
\]

直觉上，这是在版图里提前抵消一个已经被测量证明存在、而且能够重复出现的 systematic shift。

但这一步有一个非常严格的前提：

\begin{center}
\textbf{只能补偿“可重复的偏差”，不能用固定 correction 去消灭不可重复的随机 scatter。}
\end{center}

如果某一批的 wafer gradient 下一批就换了方向，那么把它写回 GDS 反而会制造新的系统误差。

\section{第五步：collision audit 要把 linewidth 也放进去}

只比较 nominal frequency spacing 还不够。一个简单的工程判据可以写成

\[
|f_i-f_j|
<
m\,
\max
\left(
\frac{f_i}{Q_{r,i}},
\frac{f_j}{Q_{r,j}}
\right),
\]

满足时先标记为 collision candidate。

这里 $m$ 是 readout / fitting / tone-placement 所需要的 margin，不是物理常数。不同 readout architecture、拟合算法和动态范围会允许不同的安全间隔。

因此一个 array-yield model 至少应该同时知道：

\begin{enumerate}
  \item design spacing；
  \item fabrication frequency scatter；
  \item loaded linewidth；
  \item readout margin；
  \item out-of-band / unusable-band 条件。
\end{enumerate}

这比单纯写“两个 resonance 距离小于 1 MHz 就算 collision”更容易迁移到不同系统。

\section{第六步：不要在同一批数据上既拟合又宣布成功}

toy model 可以这样做，因为它只是演示。但真实 fabrication 数据必须避免一个很常见的陷阱：用同一批 wafer 拟合 sensitivity，再用同一批 wafer 证明模型很好。

至少可以采用下面一种做法：

\begin{itemize}
  \item Batch A 拟合 systematic film/CD predictor；
  \item Batch B 不重新调参数，直接检查预测 residual；
  \item 只有当方向、幅度和 spatial structure 都具有 repeatability，才考虑写入下一版 GDS compensation。
\end{itemize}

如果 batch 数还很少，就不要急着叫它“process model”；更准确的名字是 working hypothesis。

\section{第七步：correlation 不是 root cause}

假设你发现
\[
R_\square(x,y)
\]
和
\[
\epsilon(x,y)
\]
看起来高度相关，也不能立刻得出“频率偏差就是方阻造成的”。

可能存在的 confounder 包括：

\begin{itemize}
  \item 膜厚和方阻本来就共变；
  \item deposition uniformity 与 lithography focus 也共享 wafer position；
  \item resonator mapping 有系统错配；
  \item temperature / package / readout calibration 具有位置相关偏差；
  \item 另一个没有记录的 process variable 同时驱动两者。
\end{itemize}

更强的证据来自 controlled split、独立 batch、as-fabricated EM model 和多种 metrology 同时闭合。

\begin{checkpointbox}
模型的目标不是给每个异常找一个“最像的原因”，而是把 hypothesis space 越缩越小，并且让下一批实验能真正区分剩余的几个解释。
\end{checkpointbox}

\section{v0.2 建议形成的四张标准图}

当真实数据开始积累后，每一批至少应该自动生成四类图：

\begin{enumerate}
  \item wafer map：$R_\square$ / thickness / CD；
  \item resonance map：$\epsilon_i$、$Q_i$、$Q_c$；
  \item predictor-vs-measurement：$\hat{\epsilon}_i$ 对 $\epsilon_i$；
  \item residual map：$r_i$ 在 wafer coordinate 上的空间结构。
\end{enumerate}

前三张告诉我们模型解释了什么，第四张告诉我们还漏了什么。

后续 v0.2 会在这个标准数据接口上再加入 plotting 和 measured-data adapter，而不是另起一套互不兼容的 notebook。

\begin{measurebox}
\textbf{一个最小可计算 fabrication-feedback loop：}
统一 ID 与单位；wafer metrology map；resonator map；明确 sensitivity source；predictor；residual；collision rule；跨 batch repeatability；只对稳定 systematic component 做 next-GDS correction。
\end{measurebox}

\section{这一章真正的出口}

读完以后，面对一批新的 KID wafer，不应该再只问“频率偏了多少”。更有用的问题是：

\begin{enumerate}
  \item 哪部分偏差能被 measured film/CD data 解释？
  \item 哪部分仍留在 residual 里？
  \item 这个 predictor 在下一批还能成立吗？
  \item collision 是 nominal spacing 不够，还是 fabrication scatter 太大？
  \item 下一版应该补偿 systematic bias，还是重新设计一个更低 sensitivity 的 geometry？
\end{enumerate}

当这些问题能够用同一套表格和脚本回答时，fabrication learning 才真正进入闭环。


\end{document}
